% Generated draft: requires mathematical and source review.
\documentclass[11pt]{article}
\usepackage[margin=25mm]{geometry}
\usepackage{fontspec}
\setmainfont{Latin Modern Roman}
\usepackage{amsmath,amssymb,mathtools,unicode-math}
\setmathfont{Latin Modern Math}
\usepackage{enumitem,needspace,xurl,hyperref}
\setlength{\parindent}{0pt}
\setlength{\parskip}{5pt}
\setlength{\emergencystretch}{3em}
\newcommand{\N}{\mathbb{N}}
\newcommand{\Z}{\mathbb{Z}}
\newcommand{\Q}{\mathbb{Q}}
\newcommand{\R}{\mathbb{R}}
\newcommand{\C}{\mathbb{C}}
\newcommand{\F}{\mathbb{F}}
\newcommand{\A}{\mathbb{A}}
\newcommand{\PP}{\mathbb{P}}
\newcommand{\E}{\mathbb{E}}
\newcommand{\eps}{\varepsilon}
\newcommand{\dd}{\,\mathrm{d}}
\begin{document}
\section*{UnsolvedMath ID: 30004097}
\textbf{Source problem code:} catalogue entry
\subsection{Definitions and mathematical statement}
Let \(k\) denote the unspecified base field, and let \(H\) be a Hopf algebra over \(k\), with coproduct \(\Delta\), counit \(\varepsilon\), and antipode \(S\). The trivial left and right \(H\)-modules \(k\) are defined through \(\varepsilon\). Under the usual ungraded convention, the Artin--Schelter Gorenstein condition means that, for some integer \(d\geq 0\),
\[
\operatorname{injdim}{}_H H=\operatorname{injdim}H_H=d,
\]
and
\[
\operatorname{Ext}^i_H(k,H)=\operatorname{Ext}^i_{H^{\mathrm{op}}}(k,H)=0\quad(i\neq d),
\qquad
\dim_k\operatorname{Ext}^d_H(k,H)=\dim_k\operatorname{Ext}^d_{H^{\mathrm{op}}}(k,H)=1.
\]
The supplied record does not specify its precise AS-Gorenstein convention or any restriction on \(k\).

An inner automorphism of \(H\) is a map \(\operatorname{Ad}_u(h)=uhu^{-1}\), where \(u\in H^\times\) is a unit. For context, a character is a unital algebra homomorphism \(\chi:H\to k\). Its left and right winding maps are commonly written, in Sweedler notation, as
\[
\tau^\ell_\chi(h)=\sum\chi(h_{(1)})h_{(2)},
\qquad
\tau^r_\chi(h)=\sum h_{(1)}\chi(h_{(2)}).
\]
These formulas explain the relevant terminology; the record does not identify \(\chi_\ell,\chi_r\) or specify which winding convention its notation \(\tau\) uses. In particular, no assignment of the displayed winding maps to the source symbols is made here.

\textbf{Requested determination.} For an AS-Gorenstein Hopf algebra \(H\) in the intended source setting, give an explicit description of the inner automorphism \(\gamma\) occurring in
\[
S^4=\gamma\circ\tau_{\chi_\ell}\circ\tau_{\chi_r}^{-1}.
\]
Equivalently, once the source characters and winding conventions have been fixed, determine a unit \(u\in H^\times\), up to multiplication by a central unit, such that
\[
S^4(h)=u\,\bigl(\tau_{\chi_\ell}\circ\tau_{\chi_r}^{-1}\bigr)(h)\,u^{-1}
\qquad\text{for every }h\in H.
\]
Here \(S^4\) denotes fourfold composition and \(\tau_{\chi_r}^{-1}\) denotes the inverse map.

\textbf{Unresolved scope.} The problem statement assumes only AS-Gorenstein, whereas the supplied background assessment concerns the noetherian AS-Gorenstein setting. Noetherianity is therefore not silently added to the statement. The hypotheses under which the displayed identity is asserted, the definitions of its characters, and the intended form of an explicit determination require source clarification.
\subsection{Short English statement}
Determine explicitly the inner conjugation in the supplied fourth-antipode formula for an Artin--Schelter Gorenstein Hopf algebra. The characters, winding conventions, and possible noetherian hypothesis remain unspecified. The supplied record labels the problem open; this draft does not independently verify that status.
\subsection{Sources}
\begin{itemize}
\item Local UnsolvedMath registry, supplied record with permanent ID \(30004097\), title ``Inner Automorphisms in the Fourth Antipode Formula,'' exact source problem code \texttt{catalogue entry}; collection: Oberwolfach Reports Open Problems. The record gives proposed year 2019 and names no proposer.
\item Supplied source citation: \emph{Mini-Workshop: Cohomology of Hopf Algebras and Tensor Categories} (2020), \emph{Oberwolfach Reports}, pp.~663--693, DOI \texttt{10.4171/owr/2019/11}. Supplied source and literature link: \url{https://doi.org/10.4171/owr/2019/11}. The record describes Oberwolfach Reports as documenting workshops and their open-problem discussions.
\item Supplied literature source: ``Report on AS-Gorenstein Hopf algebras,'' \url{https://arxiv.org/abs/2605.00774}.
\item Literature assessment dated exactly \texttt{2026-08-22}: the supplied curation labels the problem open and states that, in the general noetherian AS-Gorenstein setting, only existence of the inner automorphism is known, with an explicit description remaining open beyond special finite-dimensional situations where the distinguished grouplike gives a formula. This is supplied dated curation derived from the OWR record, not an independently verified literature assessment; its status and cited claims require independent checking.
\item The supplied registry labels the statement \texttt{corrected\_verified} and explains that notation, typography, and sentence structure were repaired while mathematical content was preserved. This is a supplied editorial qualification, not mathematical verification of this draft.
\end{itemize}
\end{document}
